\section{多边形}\label{sec:polygon}

逆时针存储的多边形：面积、凸性、点与多边形的位置关系，以及凸包、凸多边形 Minkowski 和与半平面交。

\subsection{接口}

\apitable{
    \api{Polygon<T>{p}}{顶点序列，要求逆时针；构造时自动去尾重复点并规范方向}
    \api{Polygon(vector<Vec2<T>>)}{按给定顶点构造；空序列合法}
    \api{vertices() / size() / operator[]}{只读访问顶点}
    \api{edges()}{返回每条边对应的 \code{Segment2}}
    \api{area()}{有向面积的绝对值，整数坐标返回 \code{long double}}
    \api{is_convex()}{逆时针意义下是否凸；少于 3 个点时返回 true}
    \api{relation(q)}{返回 \code{outside} / \code{boundary} / \code{inside}}
    \api{ConvexHull: convex_hull(points)}{Graham 凸包；会丢弃共线点，返回严格凸包}
    \api{minkowski_sum(a, b)}{两个凸多边形的 Minkowski 和；要求两者都凸}
    \api{half_plane_intersection(lines)}{返回 \code{optional<Polygon>}；左侧为可行域，仅浮点 \code{T}}
    \api{PointPolygonRelation}{枚举：\code{outside}、\code{boundary}、\code{inside}}
}

\subsection{约定}

顶点必须逆时针；用构造函数传入顶点时会去掉与首点重合的尾点并自动规范方向，直接写 \code{p} 字段时需自行保证顺序。\code{area()} 返回有向面积的绝对值，整数坐标返回 \code{long double}。\code{convex_hull} 丢弃共线点，返回严格凸包；\code{minkowski_sum} 只在调试构建中断言输入为凸，非凸输入不保证结果。\code{half_plane_intersection} 只接受浮点坐标，并且只会在交集为有界、非退化多边形时返回结果：空集、无界区域、点或线段都返回 \code{nullopt}。

\fullsource{geometry/polygon.hpp}
